% Chapter 1 — Theoretical foundations of external web-app penetration testing
\chapter[Theoretical Foundations]{Theoretical Foundations of External Penetration Testing of Web Applications}
\label{chap:01}
\setlength{\parskip}{1em}

This chapter establishes the theoretical foundation for the framework developed in the rest of the work. It introduces the concept of external penetration testing, surveys the standard methodology used in the industry, discusses the OWASP Top~10 catalogue of web-application risks, summarizes the dominant vulnerability classes that an external tester is expected to find, reviews the existing open-source toolchain, and finally introduces the role of Large Language Models in modern security analysis.

\section{Penetration Testing and Its External Variant}

A \textbf{penetration test} is an authorized simulated attack against an information system, performed to evaluate the system's security posture. The defining property of a penetration test, as opposed to a vulnerability scan, is that the tester actively attempts to exploit weaknesses in order to demonstrate concrete impact, not merely to enumerate potential weaknesses~\cite{nist_sp_800_115, ptes}.

Penetration tests are commonly classified by the information available to the tester:
\begin{itemize}
  \item \textbf{White-box} (or \textit{crystal-box}) testing: the tester has full access to source code, design documents, and credentials. This variant offers the deepest coverage but does not represent the view of a real external attacker.
  \item \textbf{Grey-box} testing: the tester has limited internal information, typically a normal user account and a high-level architecture description.
  \item \textbf{Black-box} testing: the tester has no internal information at all and works with only what can be observed from outside the perimeter.
\end{itemize}

\textbf{External penetration testing} corresponds to the black-box variant restricted to the public-Internet attack surface. The tester operates from a remote workstation, sees only what an unauthenticated attacker would see, and is allowed to interact only with explicitly in-scope hosts. This perspective is valuable for two reasons. First, it closely approximates the situation of a real opportunistic attacker scanning the Internet for exposed services. Second, it produces findings that are immediately actionable for the defender: every vulnerability discovered externally is, by definition, a vulnerability that any third party could also discover.

\section{Standard External-Pentest Methodology}

Industry-standard methodologies such as the \textit{Penetration Testing Execution Standard} (PTES)~\cite{ptes} and the \textit{OWASP Web Security Testing Guide} (WSTG)~\cite{owasp_wstg} describe a penetration test as a sequence of structured phases. For an external web-application engagement, the phases can be summarized as follows.

\begin{enumerate}
  \item \textbf{Pre-engagement and scoping.} The tester and the client agree in writing on the in-scope domains, IP addresses, and time window; on the level of intrusiveness (read-only versus state-changing requests); on the contact channels for incidents; and on the legal terms of engagement. The output of this phase is an unambiguous scope definition that every subsequent action must respect.
  \item \textbf{Reconnaissance (passive and active).} The tester gathers public information about the target: subdomains, DNS records, exposed ports and services, technology stack, and historical artifacts such as cached pages, public source-control repositories, and source maps. Passive reconnaissance uses third-party data sources only; active reconnaissance issues queries directly to the target (DNS resolution, HTTP requests, port scans).
  \item \textbf{Enumeration and fingerprinting.} The tester maps the live attack surface: every reachable URL, the server software behind it, the framework or CMS, the third-party libraries used by the front end, the API endpoints, the input parameters accepted by each endpoint, and the authentication mechanisms in use.
  \item \textbf{Vulnerability identification.} The tester probes the enumerated surface for known vulnerability classes (Section~\ref{sec:owasp-top10}), using a combination of automated tools, custom scripts, and manual review of responses.
  \item \textbf{Exploitation and verification.} For each candidate finding, the tester attempts a controlled exploitation that demonstrates impact without causing damage (for example, returning a benign \texttt{document.cookie} value in a confirmed reflected XSS).
  \item \textbf{Post-exploitation, if in scope.} The tester explores the consequences of a confirmed weakness: lateral movement, data extraction, or privilege escalation. For most external web-application engagements this phase is intentionally narrow, because the engagement aim is to identify weaknesses, not to compromise production data.
  \item \textbf{Reporting and remediation support.} The tester delivers a written report with executive summary, technical findings, evidence, severity ratings, and remediation recommendations, and is then available to support the client during the fix-and-retest cycle.
\end{enumerate}

The ReconX framework described in this thesis is intended to assist the analyst through phases 2--5 of this methodology. Phase~1 (scoping) and phase~7 (final reporting) remain the responsibility of the human operator, even though the framework provides explicit scope-enforcement controls and produces an HTML / JSON / Telegram report bundle that can be reused as a base for the human report.

\section{The OWASP Top~10 Catalogue of Web-Application Risks}
\label{sec:owasp-top10}

The \textbf{OWASP Top~10} is a community-driven, periodically updated list of the most critical risks to web applications~\cite{owasp_top10}. In its 2021 revision, the list contains the following categories.

\begin{description}[style=multiline, leftmargin=2.4cm, labelwidth=2cm]
  \item[A01:2021] Broken Access Control --- failures to enforce that an authenticated user can only act on resources they are allowed to access. Examples include IDOR (Insecure Direct Object References), missing function-level authorization, and cross-tenant data exposure.
  \item[A02:2021] Cryptographic Failures --- issues caused by the lack of, or the misuse of, cryptography. Examples include sensitive data transmitted in clear text, weak ciphers, and missing TLS configuration.
  \item[A03:2021] Injection --- attacker-controlled data interpreted as code by an interpreter. SQL injection, command injection, server-side template injection, and LDAP injection all fall under this category.
  \item[A04:2021] Insecure Design --- structural weaknesses that arise from the absence of secure-by-default patterns, threat modeling, or trust boundaries during the design phase.
  \item[A05:2021] Security Misconfiguration --- weak default settings, verbose error pages, unnecessary services, missing security headers, and out-of-date software.
  \item[A06:2021] Vulnerable and Outdated Components --- the use of third-party libraries or components that contain known vulnerabilities.
  \item[A07:2021] Identification and Authentication Failures --- weaknesses in the way users are identified or authenticated, including credential stuffing, weak password recovery, and broken session management.
  \item[A08:2021] Software and Data Integrity Failures --- code or infrastructure that relies on untrusted sources without integrity verification; insecure deserialization is a representative example.
  \item[A09:2021] Security Logging and Monitoring Failures --- the absence or insufficiency of logs, monitoring, and alerting required to detect, investigate, and respond to a breach.
  \item[A10:2021] Server-Side Request Forgery (SSRF) --- the ability for an attacker to make the server issue arbitrary requests to internal or external resources.
\end{description}

Each finding produced by ReconX is mapped to one of these categories through the finding-registry layer described in Chapter~\ref{chap:02}. This mapping is important because it gives every output a stable, vendor-neutral label that downstream consumers (security engineers, compliance teams, defect-tracking systems) already understand.

\section{Selected Vulnerability Classes in Detail}

The following subsections describe the vulnerability classes that the framework probes most aggressively. They are presented at the level of detail required to motivate the probes implemented in Chapter~\ref{chap:03}; a complete attack-pattern reference is outside the scope of this work.

\subsection{Reflected Cross-Site Scripting (XSS)}

\textbf{Reflected XSS} is a class of injection vulnerabilities in which the attacker delivers script content through a request parameter and the server reflects that content back into the response without proper encoding, so that the browser executes it inside the victim's session~\cite{owasp_xss, portswigger_xss}. A typical proof-of-concept request looks like

\begin{lstlisting}[basicstyle=\ttfamily\small]
GET /search?q=<script>alert(1)</script> HTTP/1.1
\end{lstlisting}

For a black-box scanner the practical task is not to prove a complete \texttt{alert(1)} exploitation chain, but to determine whether a unique attacker-controlled marker is reflected into the response in a context that would allow script execution. ReconX uses a unique per-scan marker (\texttt{rxss}\textit{<random hex>}) embedded inside several breakout payloads (plain reflection, HTML attribute breakout, script-tag breakout, single-quote breakout) and inspects the response for the marker. When a credible dynamic-analysis tool such as \texttt{dalfox}~\cite{dalfox} is available on the operator's machine, ReconX additionally invokes it for parameters that look promising.

\subsection{SQL Injection (SQLi)}

\textbf{SQL injection} occurs when user-controlled input is concatenated into a SQL query string without parameterization, allowing the attacker to alter the structure of the query~\cite{owasp_sqli}. The standard external-pentest tool for SQL-injection identification is \texttt{sqlmap}~\cite{sqlmap}, a research-grade probe that supports boolean-based blind, time-based blind, error-based, UNION-based, stacked-query, and inline-query techniques across a wide list of database management systems. Because \texttt{sqlmap} is itself an active and noisy probe, ReconX wraps it in a conservative configuration (level~1, risk~1, smart detection, random user agent) and limits the number of targets per scan to a small upper bound (default~10).

\subsection{Insecure Direct Object Reference (IDOR) and Broken Object-Level Authorization}

An \textbf{IDOR} vulnerability exists when a request that operates on a specific object trusts an attacker-controlled identifier without verifying that the authenticated user is authorized to act on that object~\cite{owasp_bola}. Because the underlying authorization gap is the issue, IDOR cannot be detected purely by inspecting one user's view of the application; the detector must compare two distinct views and verify that one of them returns data that should be visible only to the other.

ReconX handles this in two complementary ways. First, it produces \textit{candidate} findings by enumerating request parameters whose names suggest an object identifier (\texttt{user\_id}, \texttt{account\_id}, \texttt{order\_id}, \texttt{uuid}, integer-only parameters, and similar patterns) and recording them for manual review. Second, when the operator has configured an \texttt{auth\_profiles} block with two or more profiles (a low-privilege user and an admin), ReconX performs cross-profile comparison: it fetches each candidate URL once as user~A, once as user~B, and once anonymously, and computes a textual similarity score between the responses. A surprisingly high cross-profile similarity is reported as a confirmed candidate; this is the same logic that an experienced human tester uses.

\subsection{Host Header Injection and Web Cache Poisoning}

The \textbf{Host} HTTP header is a primary source of trust when the server constructs absolute URLs, builds password-reset links, or routes requests behind a reverse proxy~\cite{portswigger_host}. Web frameworks that reflect the Host header (or \texttt{X-Forwarded-Host}, \texttt{Forwarded}, \texttt{X-Original-Host}) without validation can be tricked into emitting attacker-controlled hostnames into responses, which an attacker can then weaponize against other users via cached responses (a class known as \textbf{web cache poisoning}~\cite{portswigger_cache}).

ReconX issues each request with a battery of host-related headers carrying a unique marker, looks for the marker in the response body or in the \texttt{Location} header, and inspects \texttt{Age}, \texttt{X-Cache}, and similar headers to assess whether the response was returned from a cache. Severity is escalated when reflection and cacheability co-occur, because the combination indicates a likely cache-poisoning chain.

\subsection{Server-Side Request Forgery (SSRF) and Server-Side Template Injection (SSTI)}

\textbf{SSRF} is a vulnerability in which the server fetches a URL controlled by the attacker, allowing the attacker to reach internal resources from the server's privileged network position~\cite{owasp_ssrf}. A common SSRF payload targets cloud-instance metadata endpoints such as \texttt{169.254.169.254/latest/meta-data/}, which return credentials usable to pivot inside the cloud account.

\textbf{SSTI} is a class of injection vulnerabilities in which user input is interpreted by a server-side template engine (Jinja2, Twig, Freemarker, ERB), allowing the attacker to execute template code on the server~\cite{portswigger_ssti}. A canonical probe is to submit \verb|{{7*7}}| and observe whether the server returns \texttt{49}; further payloads escalate from arithmetic to function calls.

ReconX implements both classes in the \texttt{injection\_probe} module: a small fixed set of SSTI probes covering the most common template syntaxes, and a set of SSRF probes that target the cloud-metadata endpoints and the most exposed loopback ports (SSH, Redis, MongoDB). Out-of-band confirmation of SSRF is delegated to Nuclei's \texttt{interactsh} integration when the operator has enabled it.

\subsection{XML External Entity (XXE), Deserialization, and HTTP Request Smuggling}

\textbf{XXE} occurs when an XML parser dereferences external entities controlled by the attacker, typically to exfiltrate local files or to perform a server-side request~\cite{owasp_xxe}. \textbf{Insecure deserialization} occurs when an application deserializes attacker-controlled data without integrity checks, often resulting in remote code execution~\cite{owasp_deserialization}. \textbf{HTTP request smuggling} arises when a chain of front-end and back-end servers disagrees about how to delimit successive HTTP requests on the same TCP connection, allowing the attacker to inject a request that is parsed differently by each component~\cite{portswigger_smuggling}.

These three classes are difficult to identify reliably from the outside, because a non-destructive probe can usually only collect indirect evidence (parsing-error signatures, response-timing deviations, deserialization stack-trace fragments). ReconX therefore reports them with low to medium confidence by default, and delegates aggressive confirmation to the operator's manual follow-up.

\subsection{Other Configuration-Level Issues}

In addition to the active probes listed above, an external scanner is expected to surface several configuration-level issues whose existence can be determined from a single request, without active payload crafting. These include weak or expired TLS certificates, missing security headers (\texttt{Strict-Transport-Security}, \texttt{Content-Security-Policy}, \texttt{X-Content-Type-Options}, \texttt{X-Frame-Options}), permissive CORS policies (wildcard or reflected \texttt{Access-Control-Allow-Origin} with \texttt{Allow-Credentials: true}), unsigned or weakly signed JWT tokens, hardcoded secrets in publicly served JavaScript files, and subdomains pointing to unclaimed third-party services (subdomain takeover).

\section{Existing Open-Source Toolchain}
\label{sec:toolchain}

Modern external web-application penetration testing is supported by a large ecosystem of open-source tools, each typically optimized for a single phase or vulnerability class. A non-exhaustive overview of the tools relevant to this work is given in Table~\ref{tab:tools}.

\begin{table}[h]
\centering
\footnotesize
\begin{tabular}{|l|l|l|}
\hline
\textbf{Tool} & \textbf{Phase / class} & \textbf{Role in ReconX} \\
\hline
\texttt{subfinder} \cite{subfinder}     & passive subdomain enumeration  & invoked by recon module \\
\texttt{amass} \cite{amass}              & subdomain enumeration + DNS    & invoked by recon module \\
\texttt{assetfinder}                     & subdomain enumeration          & invoked by recon module \\
\texttt{nmap} \cite{nmap}                & port and service scanning      & invoked by portscan module \\
\texttt{masscan}                         & high-rate port scanning        & alternative to nmap \\
\texttt{wappalyzer} \cite{wappalyzer}    & technology fingerprinting      & embedded JSON database \\
\texttt{wpscan} \cite{wpscan}            & WordPress vulnerability scan   & invoked by cmscan module \\
\texttt{katana}                          & web crawling                   & invoked by fuzzer module \\
\texttt{feroxbuster} / \texttt{ffuf}     & directory brute-forcing        & invoked by fuzzer module \\
\texttt{arjun} \cite{arjun}              & hidden-parameter discovery     & invoked by parameter\_discovery \\
\texttt{dalfox} \cite{dalfox}            & dynamic XSS testing            & invoked by xss module \\
\texttt{sqlmap} \cite{sqlmap}            & SQL injection testing          & invoked by sql\_injection module \\
\texttt{nuclei} \cite{nuclei}            & template-based vulnerability scan & invoked by vulnscan module \\
\texttt{gitleaks}                        & secret scanning in git history & invoked by secret\_scanner \\
\texttt{searchsploit}                    & ExploitDB lookup               & invoked by cve\_check \\
\texttt{interactsh} \cite{interactsh}    & out-of-band callback server    & optional Nuclei dependency \\
\hline
\end{tabular}
\caption{Selected open-source tools used by ReconX}
\label{tab:tools}
\end{table}

A defining feature of this ecosystem is its fragmentation: each tool produces its own output format, its own concept of severity, its own concept of evidence, and its own command-line interface. The analyst is left to glue them together manually for every engagement. One of the contributions of the present work is to provide that glue --- not in the form of a thin wrapper, but in the form of a deeper integration in which every tool's output is parsed, normalized, deduplicated, and finally combined into a single ranked finding list.

\section{The Role of Large Language Models}

The recent availability of capable open-weight Large Language Models (LLMs) such as DeepSeek-R1, Qwen~2.5, and Llama~3 changes what is reasonable to expect from a security tool's reporting layer~\cite{ollama, deepseek_r1}. A few years ago a security tool could only emit raw findings, and the analyst had to invest hours of writing to turn them into an executive summary. With a locally running LLM, the same tool can now generate, in seconds, a coherent natural-language analysis that ranks the findings, explains why a particular combination matters, and proposes remediation steps in a tone appropriate for the intended audience (executive, engineer, or auditor).

LLMs do not replace the security analyst. They are not authoritative about exploitability, they do not know what is in scope, and they may hallucinate. But they substantially accelerate the steps that consist of \textit{rewriting structured data as readable prose}, which is most of the writing burden of a security report. In ReconX the LLM is used in two places:

\begin{enumerate}
  \item \textbf{Real-time advice} during a long scan: after each major module completes, the framework can send a short structured summary to a local LLM and post the model's three to six bullet advice via Telegram. This keeps the analyst engaged and lets them adjust scan parameters mid-flight.
  \item \textbf{Post-scan analysis}: at the end of the scan, the framework aggregates the top findings across every module and asks the LLM to produce a Markdown analytical report with an executive summary, risk tiers, and prioritized investigation steps. The result is embedded into the final HTML report.
\end{enumerate}

\section{Research Gap and Rationale for the Selected Approach}

The literature and the existing toolchain reviewed in this chapter motivate the design choices in the chapters that follow. Specifically:

\begin{itemize}
  \item No single open-source tool covers the breadth of phases (reconnaissance, configuration analysis, active probing, correlation, reporting) needed for an external pentest. Existing frameworks either focus narrowly on one phase or aggregate other tools without a unified finding model.
  \item There is value in a framework that maintains a single, normalized finding schema across all phases, so that downstream consumers (correlator, asset-risk module, report renderer, AI report module) do not need to know about each probe's output format.
  \item There is value in a framework that enforces scope at the lowest possible layer, so that no probe --- including third-party probes invoked via subprocess --- can accidentally send a request to an out-of-scope host.
  \item There is value in an explicit \textit{verdict} layer that separates \textit{confirmed} from \textit{candidate} findings, because the cost of a false positive in a security report (analyst time, reputational damage with the client) is far higher than the cost of a missed weak signal.
  \item Local LLMs are now mature enough to be integrated as a first-class reporting component, but the integration must be designed so that the framework still produces a usable report when the LLM is unavailable.
\end{itemize}

The next chapter translates these observations into a concrete framework design.
